%% toc: P1 4 4.5,4.6
\section{Proofs for finite coding over an infinite cardinal and syntax counting}
\label{sec:annex-p01-w02-p01-011-finite-coding}

Notation and readings are those of \cref{def:p01-w02-p01-011-readings}. The
Cantor--Schr\"oder--Bernstein theorem of \cref{sec:set-prerequisites-05} is
cited as CSB; its proof there uses no choice. Equality of ordered pairs is
the ordered-pair identity of \cref{sec:set-prerequisites-02}. A strict
well-order restricted to a subset is again a strict well-order:
irreflexivity, transitivity and comparability are statements about members,
and a nonempty subset of the subset is a nonempty subset of the whole set.
This remark is used without further mention.

\begin{annexitem}{annex:p01-w02-p01-011-ordinals}{thm:p01-w02-p01-011-ordinals}{Basic facts about ordinals}
\begin{proof}
\emph{Item~1.} Let \(x\in\alpha\). Since \(\alpha\) is transitive,
\(x\subseteq\alpha\), so \(\in\) strictly well-orders \(x\). To see that
\(x\) is transitive, let \(z\in y\in x\). Then \(y\in\alpha\) because
\(x\subseteq\alpha\), and \(z\in\alpha\) because \(\alpha\) is transitive.
Thus \(x,y,z\in\alpha\) with \(z\in y\) and \(y\in x\), and transitivity of
the relation \(\in\) on \(\alpha\) gives \(z\in x\). So \(x\) is an ordinal.
If \(\alpha\in\alpha\), then \(\alpha\) would be a member of \(\alpha\)
standing in the relation \(\in\) to itself, against irreflexivity of \(\in\)
on \(\alpha\).

\emph{Item~2.} Let \(\alpha\subseteq\beta\) and \(\alpha\ne\beta\). The set
\(\beta\setminus\alpha\) is a nonempty subset of \(\beta\); let \(\gamma\) be
its \(\in\)-least member. We show \(\gamma=\alpha\). Let \(\delta\in\gamma\).
Then \(\delta\in\beta\) since \(\beta\) is transitive; if \(\delta\notin\alpha\),
then \(\delta\in\beta\setminus\alpha\) and \(\delta\in\gamma\), against the
choice of \(\gamma\). So \(\gamma\subseteq\alpha\). Conversely let
\(\delta\in\alpha\). Then \(\delta,\gamma\in\beta\), so \(\delta\in\gamma\),
\(\delta=\gamma\) or \(\gamma\in\delta\). The middle case is impossible
because \(\delta\in\alpha\) and \(\gamma\notin\alpha\). In the last case
\(\gamma\in\alpha\) because \(\alpha\) is transitive and \(\delta\in\alpha\),
which is again impossible. So \(\delta\in\gamma\), and
\(\alpha\subseteq\gamma\). Hence \(\alpha=\gamma\in\beta\).

\emph{Item~3.} Put \(\iota=\alpha\cap\beta\). If \(y\in x\in\iota\), then
\(y\in\alpha\) and \(y\in\beta\) because both are transitive; so \(\iota\) is
transitive. As \(\iota\subseteq\alpha\), \(\in\) strictly well-orders
\(\iota\). Hence \(\iota\) is an ordinal. Item~2 applied to
\(\iota\subseteq\alpha\) and to \(\iota\subseteq\beta\) gives
\((\iota=\alpha\text{ or }\iota\in\alpha)\) and
\((\iota=\beta\text{ or }\iota\in\beta)\). If \(\iota\in\alpha\) and
\(\iota\in\beta\), then \(\iota\in\iota\), against item~1 for \(\iota\). So
\(\iota=\alpha\) or \(\iota=\beta\). If \(\iota=\alpha\), then
\(\alpha\subseteq\beta\) and item~2 gives \(\alpha=\beta\) or
\(\alpha\in\beta\). If \(\iota=\beta\), then \(\beta\subseteq\alpha\) and
\(\beta=\alpha\) or \(\beta\in\alpha\). So at least one of the three cases
holds. At most one holds: \(\alpha\in\beta\) together with \(\alpha=\beta\)
gives \(\alpha\in\alpha\); \(\beta\in\alpha\) together with \(\alpha=\beta\)
gives the same; and \(\alpha\in\beta\) together with \(\beta\in\alpha\) gives
\(\beta\subseteq\alpha\) by transitivity of \(\alpha\), hence
\(\alpha\in\alpha\). Each contradicts item~1.

For the last claim, if \(\alpha=\beta\) then \(\alpha\subseteq\beta\), and if
\(\alpha\in\beta\) then \(\alpha\subseteq\beta\) because \(\beta\) is
transitive. Conversely, \(\alpha\subseteq\beta\) gives \(\alpha\leq\beta\) by
item~2.

\emph{Item~4.} Let \(\alpha\in S\). If \(\alpha\cap S=\varnothing\), then no
\(\delta\in S\) lies in \(\alpha\), and item~3 gives \(\alpha=\delta\) or
\(\alpha\in\delta\), that is \(\alpha\leq\delta\), for every \(\delta\in S\).
Otherwise \(\alpha\cap S\) is a nonempty subset of \(\alpha\); let
\(\gamma\) be its \(\in\)-least member. Let \(\delta\in S\). If
\(\delta\in\alpha\), then \(\delta\in\alpha\cap S\), so \(\delta\in\gamma\)
fails, and item~3 gives \(\gamma\leq\delta\). If \(\delta\notin\alpha\), then
item~3 gives \(\alpha\leq\delta\), so \(\alpha\subseteq\delta\), and
\(\gamma\in\alpha\) gives \(\gamma\in\delta\). In both cases
\(\gamma\leq\delta\). If \(\gamma\) and \(\gamma'\) are both least, then
\(\gamma\leq\gamma'\) and \(\gamma'\leq\gamma\); if they were distinct this
would give \(\gamma\in\gamma'\) and \(\gamma'\in\gamma\), which item~3
excludes.

\emph{Item~5.} Let \(T\) be a transitive set of ordinals. The relation \(\in\)
on \(T\) is irreflexive by item~1. It is transitive: if \(x\in y\in z\) with
\(x,y,z\in T\), then \(x\in z\) because the ordinal \(z\) is transitive. Any
two distinct members are comparable by item~3. A nonempty \(S\subseteq T\)
has a least member \(\gamma\) by item~4, and no \(\delta\in S\) satisfies
\(\delta\in\gamma\), since \(\gamma\leq\delta\) together with
\(\delta\in\gamma\) is excluded by item~3. So \(\in\) strictly well-orders
\(T\), and \(T\) is an ordinal. The members of \(\alpha+1=\alpha\cup\{\alpha\}\)
are the members of \(\alpha\), which are ordinals by item~1, and \(\alpha\).
The set is transitive: if \(y\in x\in\alpha+1\), then either \(x\in\alpha\)
and \(y\in\alpha\) because \(\alpha\) is transitive, or \(x=\alpha\) and
\(y\in\alpha\). In both cases \(y\in\alpha+1\). So \(\alpha+1\) is an
ordinal, and \(\gamma\in\alpha+1\) iff \(\gamma\in\alpha\) or
\(\gamma=\alpha\), that is iff \(\gamma\leq\alpha\).

\emph{Item~6.} By induction on \(n\). The set \(0=\varnothing\) is transitive
and is strictly well-ordered by \(\in\) vacuously. If \(n\) is an ordinal,
then \(n+1\) is an ordinal by item~5. Hence \(n\notin n\) by item~1. Let
\(m+1=n+1\) and suppose \(m\ne n\). From \(m\in n+1\) and \(m\ne n\) we get
\(m\in n\); from \(n\in m+1\) and \(n\ne m\) we get \(n\in m\). This
contradicts item~3 for the ordinals \(m\) and \(n\).
\end{proof}
\end{annexitem}

\begin{annexitem}{annex:p01-w02-p01-011-finite-sets}{thm:p01-w02-p01-011-finite-sets}{Finite sets}
\begin{proof}
\emph{Item~1.} Let \(c:r\to C\) be a bijection. Since \(r\notin r\)
(item~6 of \cref{thm:p01-w02-p01-011-ordinals}),
\(c'=c\cup\{(r,x)\}\) is a function with domain \(r\cup\{r\}=r+1\). Its
range is \(C\cup\{x\}\), and it is injective because \(c\) is injective and
\(x\notin C=c[r]\).

\emph{Item~2.} We prove by induction on \(q\in\N\): for every finite \(A\)
and every \(B\) with \(\abs B=\abs q\) and \(A\cap B=\varnothing\), the set
\(A\cup B\) is finite. For \(q=0\), \(B\) is the range of the empty function,
so \(B=\varnothing\) and \(A\cup B=A\). Assume the claim for \(q\), and let
\(b:q+1\to B\) be a bijection. Put \(B'=b[q]\) and \(y=b(q)\). The
restriction of \(b\) to \(q\) is a bijection \(q\to B'\); \(y\notin B'\)
because \(b\) is injective and \(q\notin q\); and \(B=B'\cup\{y\}\). Since
\(A\cap B'=\varnothing\), the induction hypothesis gives
\(\abs{A\cup B'}=\abs r\) for some \(r\in\N\). As \(y\in B\) and
\(A\cap B=\varnothing\), \(y\notin A\); so \(y\notin A\cup B'\), and item~1
gives \(\abs{A\cup B}=\abs{(A\cup B')\cup\{y\}}=\abs{r+1}\).

\emph{Item~3.} First, by induction on \(n\), every \(S\subseteq n\) is
finite. For \(n=0\), \(S=\varnothing=0\). Let \(S\subseteq n+1\). The set
\(S\cap n\subseteq n\) is finite by the induction hypothesis, say
\(\abs{S\cap n}=\abs r\). If \(n\notin S\), then \(S\subseteq n\) and
\(S=S\cap n\). If \(n\in S\), then \(S=(S\cap n)\cup\{n\}\) with
\(n\notin S\cap n\) because \(n\notin n\), and item~1 gives
\(\abs S=\abs{r+1}\). Now let \(e:n\to A\) be a bijection and
\(S\subseteq A\). Then \(T=e^{-1}[S]\subseteq n\) is finite, say via a
bijection \(d:r\to T\), and \(e\circ d:r\to S\) is a bijection because
\(e\) restricted to \(T\) is a bijection \(T\to S\).

\emph{Item~4.} By induction on \(k\). For \(k=0\) the union is empty. Let
\(A_i\) be finite for every \(i<k+1\), and put
\(U=\bigcup_{i<k}A_i\), finite by the induction hypothesis applied to the
family restricted to \(k\subseteq k+1\). Then
\(\bigcup_{i<k+1}A_i=U\cup(A_k\setminus U)\). The set \(A_k\setminus U\) is
finite by item~3 and is disjoint from \(U\), so item~2 applies.

\emph{Item~5.} Let \(a:n\to A\) and \(b:m\to B\) be bijections. For
\(i<n\), the map \(j\mapsto(a(i),b(j))\) is a bijection
\(m\to\{a(i)\}\times B\), because \(b\) is a bijection and pairs are equal
only when their components are equal. Since \(a\) is onto \(A\),
\(A\times B=\bigcup_{i<n}\{a(i)\}\times B\), which is finite by item~4.
\end{proof}
\end{annexitem}

\begin{annexitem}{annex:p01-w02-p01-011-infinite-ordinals}{thm:p01-w02-p01-011-infinite-ordinals}{Infinite ordinals and cardinals}
\begin{proof}
\emph{Item~1.} We prove \(n\in\alpha\) by induction on \(n\in\N\). The ordinal
\(\alpha\) is not empty, since \(\varnothing=0\) is bijective with
\(0\in\N\). Let \(\xi\) be the least member of \(\alpha\)
(item~4 of \cref{thm:p01-w02-p01-011-ordinals}). If \(\eta\in\xi\), then
\(\eta\in\alpha\) because \(\alpha\) is transitive, and \(\eta<\xi\)
contradicts the choice of \(\xi\). So \(\xi=\varnothing=0\), and
\(0\in\alpha\). Now let \(n\in\alpha\). Then \(n\subseteq\alpha\) because
\(\alpha\) is transitive, so \(n+1=n\cup\{n\}\subseteq\alpha\). The set
\(n+1\) is an ordinal (item~6 of \cref{thm:p01-w02-p01-011-ordinals}), and
\(n+1\ne\alpha\) because \(\alpha\) is bijective with no natural number. By
item~2 of \cref{thm:p01-w02-p01-011-ordinals}, \(n+1\in\alpha\). Hence
\(\N\subseteq\alpha\). If \(X\) is countable, there is an injection
\(X\to\N\), and composing it with the inclusion \(\N\subseteq\alpha\) gives
an injection \(X\to\alpha\). For \(\alpha=\kappa\) and \(X=\N\), the inclusion
itself is an injection \(\N\to\kappa\), which is \(\aleph_0\leq\kappa\).

\emph{Item~2.} By item~1, \(\N\subseteq\alpha\), and \(\alpha\notin\N\)
because \(\alpha\notin\alpha\). So \(\alpha+1\) is the disjoint union of
\(\{\alpha\}\), \(\N\) and \(\alpha\setminus\N\). Define
\(f:\alpha+1\to\alpha\) by \(f(\alpha)=0\), \(f(n)=n+1\) for \(n\in\N\),
and \(f(\xi)=\xi\) for \(\xi\in\alpha\setminus\N\). Its values lie in
\(\alpha\) because \(\N\subseteq\alpha\). On \(\{\alpha\}\cup\N\) the values
lie in \(\N\) and are pairwise distinct: \(n+1\ne0\) since \(n\in n+1\), and
\(m+1=n+1\) implies \(m=n\) (item~6 of
\cref{thm:p01-w02-p01-011-ordinals}). On \(\alpha\setminus\N\), \(f\) is the
identity, with values outside \(\N\). So \(f\) is injective. It is
surjective: the set \(\{0\}\cup\{m+1:m\in\N\}\) contains \(0\) and is closed
under \(n\mapsto n+1\), so it is all of \(\N\); thus every member of \(\N\)
is \(0=f(\alpha)\) or \(m+1=f(m)\), and every
\(\xi\in\alpha\setminus\N\) is \(f(\xi)\).

\emph{Item~3.} Let \(\gamma\in\kappa\). Then \(\gamma\) is an ordinal,
\(\gamma\subseteq\kappa\) and so \(\gamma+1\subseteq\kappa\); and
\(\gamma+1\) is an ordinal. Suppose \(\gamma+1=\kappa\). If \(\gamma\) were
bijective with some \(n\in\N\), item~1 of
\cref{thm:p01-w02-p01-011-finite-sets} (with \(\gamma\notin\gamma\)) would
make \(\kappa=\gamma+1\) bijective with \(n+1\in\N\), against the
infiniteness of \(\kappa\). So \(\gamma\) is infinite, and item~2 gives
\(\abs\kappa=\abs{\gamma+1}=\abs\gamma\) with \(\gamma<\kappa\), against the
definition of a cardinal. Hence \(\gamma+1\ne\kappa\), and item~2 of
\cref{thm:p01-w02-p01-011-ordinals} gives \(\gamma+1\in\kappa\).

\emph{Item~4.} The set \(S=\{\xi\in\beta+1:\abs\xi=\abs\beta\}\) exists by
separation and contains \(\beta\). Let \(\operatorname{card}(\beta)\) be its least member.
Then \(\operatorname{card}(\beta)\leq\beta\) by item~5 of
\cref{thm:p01-w02-p01-011-ordinals}. If an ordinal \(\zeta<\operatorname{card}(\beta)\) were
bijective with \(\beta\), then \(\zeta\in \operatorname{card}(\beta)\subseteq\beta+1\) would
put \(\zeta\) in \(S\) below its least member; so \(\operatorname{card}(\beta)\) is the least
ordinal bijective with \(\beta\). It is a cardinal: an ordinal
\(\zeta<\operatorname{card}(\beta)\) bijective with \(\operatorname{card}(\beta)\) would be bijective with
\(\beta\). Finally, a bijection \(\operatorname{card}(\beta)\to\beta\) composed with a
bijection between \(\beta\) and some \(n\in\N\) is a bijection between
\(\operatorname{card}(\beta)\) and \(n\), and conversely; so \(\operatorname{card}(\beta)\) is infinite iff
\(\beta\) is.

\emph{Item~5.} Let \(g:X\to\beta\) be an injection. Since \(\beta\in\kappa\)
and \(\kappa\) is transitive, \(\beta\subseteq\kappa\), so \(g\) is also an
injection \(X\to\kappa\). Suppose \(h:X\to\kappa\) were a bijection. Then
\(g\circ h^{-1}\) is an injection \(\kappa\to\beta\), and the inclusion is
an injection \(\beta\to\kappa\); CSB gives a bijection between \(\beta\) and
\(\kappa\) with \(\beta<\kappa\), against the definition of a cardinal.

\emph{Item~6.} Let \(X\) be bijective with \(n\in\N\). By item~1,
\(n\in\kappa\), so \(n<\kappa\), and \(\abs X\leq\abs n\). Item~5 gives
\(\abs X<\kappa\).

\emph{Item~7.} Let \(e:X\to\mu\) be a bijection. If \(\mu<\kappa\), then
\(\abs X\leq\abs\mu\) through \(e\), and item~5 with \(\beta=\mu\) gives
\(\abs X<\kappa\). Conversely, let \(g:X\to\kappa\) be an injection and
suppose no bijection \(X\to\kappa\) exists. By item~3 of
\cref{thm:p01-w02-p01-011-ordinals}, \(\mu\in\kappa\), \(\mu=\kappa\) or
\(\kappa\in\mu\). If \(\mu=\kappa\), \(e\) is a bijection \(X\to\kappa\),
which is excluded. If \(\kappa\in\mu\), then \(\kappa\subseteq\mu\), so
\(e^{-1}\) restricted to \(\kappa\) is an injection \(\kappa\to X\); with
\(g\), CSB gives a bijection \(X\to\kappa\), which is excluded. So
\(\mu\in\kappa\).
\end{proof}
\end{annexitem}

\begin{annexitem}{annex:p01-w02-p01-011-canonical-order}{thm:p01-w02-p01-011-canonical-order}{The canonical order}
\begin{proof}
For ordinals \(\alpha,\beta\), item~3 of
\cref{thm:p01-w02-p01-011-ordinals} gives \(\alpha\leq\beta\) or
\(\beta<\alpha\), so \(\max(\alpha,\beta)\) is one of \(\alpha,\beta\) and
satisfies \(\alpha\leq\max(\alpha,\beta)\) and
\(\beta\leq\max(\alpha,\beta)\). The relation \(\leq\) on ordinals is
transitive, because it is \(\subseteq\) by the same item.

\emph{Item~1.} For triples of ordinals put
\(a<_{\mathrm{lex}}b\) iff there is \(i\in\{1,2,3\}\) with \(a_j=b_j\) for
every \(j<i\) and \(a_i<b_i\); such an \(i\) is the first index at which
\(a\) and \(b\) differ. The three clauses of
\cref{def:p01-w02-p01-011-canonical-order} are the cases \(i=1,2,3\) for
\(k(p)\) and \(k(q)\), so \(p\lhd q\) iff \(k(p)<_{\mathrm{lex}}k(q)\).

\(<_{\mathrm{lex}}\) is irreflexive, since \(a<_{\mathrm{lex}}a\) would need
\(a_i\in a_i\). It is transitive: let \(a<_{\mathrm{lex}}b\) at index \(i\)
and \(b<_{\mathrm{lex}}c\) at index \(j\), and let \(k=\min(i,j)\). Below
\(k\), the three triples agree. At \(k\): if \(i<j\), then
\(a_k<b_k=c_k\); if \(j<i\), then \(a_k=b_k<c_k\); if \(i=j\), then
\(a_k\in b_k\in c_k\), so \(a_k\in c_k\) because \(c_k\) is transitive. So
\(a<_{\mathrm{lex}}c\) at index \(k\). Distinct triples are comparable: at
the first index where they differ, item~3 of
\cref{thm:p01-w02-p01-011-ordinals} decides. Every nonempty set \(T\) of
triples has a least member: let \(m_1\) be the least first coordinate of
members of \(T\) (item~4 of \cref{thm:p01-w02-p01-011-ordinals}), \(m_2\) the
least second coordinate of members of \(T\) with first coordinate \(m_1\),
and \(m_3\) the least third coordinate of members of \(T\) with first two
coordinates \(m_1,m_2\). Then \(t=(m_1,m_2,m_3)\in T\), and for every other
\(t'\in T\) the first index at which \(t'\) differs from \(t\) carries a
coordinate of \(t'\) larger than that of \(t\), so
\(t<_{\mathrm{lex}}t'\); by irreflexivity and transitivity
\(t'<_{\mathrm{lex}}t\) fails.

The map \(k\) is injective on \(\mu\times\mu\), because \(k(\alpha,\beta)\)
has \(\alpha\) and \(\beta\) as its second and third coordinates. Hence
\(\lhd\) is irreflexive and transitive, distinct pairs are comparable
because their images under \(k\) are distinct, and a nonempty
\(S\subseteq\mu\times\mu\) has as \(\lhd\)-least member the pair whose image
is the least member of \(k[S]\). So \(\lhd\) strictly well-orders
\(\mu\times\mu\).

\emph{Item~2.} Let \((\xi,\eta)\lhd(\alpha,\beta)\) and
\(\gamma=\max(\alpha,\beta)\). Every clause of
\cref{def:p01-w02-p01-011-canonical-order} gives
\(\max(\xi,\eta)<\gamma\) or \(\max(\xi,\eta)=\gamma\), so
\(\max(\xi,\eta)\leq\gamma\). Then \(\xi\leq\max(\xi,\eta)\leq\gamma\) and
\(\eta\leq\max(\xi,\eta)\leq\gamma\), and item~5 of
\cref{thm:p01-w02-p01-011-ordinals} gives \(\xi,\eta\in\gamma+1\).
\end{proof}
\end{annexitem}

\begin{annexitem}{annex:p01-w02-p01-011-representation}{thm:p01-w02-p01-011-representation}{Initial segments are represented by ordinals}
\begin{proof}
\emph{Item~1.} Suppose \(f\) is not the identity. Let \(\xi\) be the least
member of \(\{\xi'\in\alpha:f(\xi')\ne\xi'\}\) (item~4 of
\cref{thm:p01-w02-p01-011-ordinals}). We show
\(f(\xi)=\{f(\eta):\eta\in\xi\}\). If \(\eta\in\xi\), then
\(f(\eta)\in f(\xi)\). Conversely let \(\zeta\in f(\xi)\). Then
\(\zeta\in\beta\) because \(\beta\) is transitive, so \(\zeta=f(\eta)\) for
some \(\eta\in\alpha\). If \(\eta=\xi\), then \(f(\xi)\in f(\xi)\), against
item~1 of \cref{thm:p01-w02-p01-011-ordinals}. If \(\xi\in\eta\), then
\(f(\xi)\in f(\eta)=\zeta\in f(\xi)\), and transitivity of the ordinal
\(f(\xi)\) gives \(f(\xi)\in f(\xi)\) again. So \(\eta\in\xi\) by item~3 of
\cref{thm:p01-w02-p01-011-ordinals}. By the choice of \(\xi\),
\(f(\eta)=\eta\) for every \(\eta\in\xi\), so
\(f(\xi)=\{\eta:\eta\in\xi\}=\xi\), against the choice of \(\xi\). Hence
\(f\) is the identity and \(\beta=f[\alpha]=\alpha\).

Consequently a strictly ordered set is order isomorphic to at most one
ordinal: if \(g\) and \(g'\) are order isomorphisms from it onto \(\alpha\)
and \(\beta\), then \(g'\circ g^{-1}\) is an order isomorphism
\(\alpha\to\beta\), so \(\alpha=\beta\).

\emph{Item~2, a restriction step.} Let \(w\in W\), let \(\theta\) be an
ordinal, let \(h:\operatorname{seg}(w)\to\theta\) be an order isomorphism,
and let \(v\prec w\). Then \(h(v)\) is an ordinal and the restriction of
\(h\) to \(\operatorname{seg}(v)\) is an order isomorphism
\(\operatorname{seg}(v)\to h(v)\). Indeed
\(\operatorname{seg}(v)\subseteq\operatorname{seg}(w)\) by transitivity of
\(\prec\), and \(u\prec v\) gives \(h(u)\in h(v)\). If \(\zeta\in h(v)\),
then \(\zeta\in\theta\), so \(\zeta=h(u)\) for some
\(u\in\operatorname{seg}(w)\); \(u=v\) would give \(h(v)\in h(v)\), and
\(v\prec u\) would give \(h(v)\in h(u)=\zeta\in h(v)\), hence
\(h(v)\in h(v)\) by transitivity of \(h(v)\); so \(u\prec v\). Thus the
restriction maps \(\operatorname{seg}(v)\) onto \(h(v)\), and as a
restriction of an order isomorphism it is injective and preserves and
reflects the order.

\emph{Item~2, existence.} Let \(D\) be the set of \(w\in W\) such that
\(\operatorname{seg}(w)\) is order isomorphic to some ordinal, and suppose
\(W\setminus D\ne\varnothing\); let \(w_0\) be its \(\prec\)-least member.
Every \(v\prec w_0\) lies in \(D\), and by item~1 there is exactly one
ordinal \(F(v)\) order isomorphic to \(\operatorname{seg}(v)\). By
replacement, \(\theta=\{F(v):v\prec w_0\}\) is a set. Let
\(u\prec v\prec w_0\), and let \(h_v:\operatorname{seg}(v)\to F(v)\) be an
order isomorphism. The restriction step makes \(\operatorname{seg}(u)\)
order isomorphic to \(h_v(u)\), so \(F(u)=h_v(u)\in F(v)\) by uniqueness.
The set \(\theta\) is transitive: if \(\zeta\in F(v)\) with \(v\prec w_0\),
take \(h_v\) as before; \(\zeta=h_v(u)\) for some \(u\prec v\), and
\(\zeta=h_v(u)=F(u)\in\theta\) because \(u\prec w_0\). So \(\theta\) is a
transitive set of ordinals, an ordinal by item~5 of
\cref{thm:p01-w02-p01-011-ordinals}. The map \(v\mapsto F(v)\) sends
\(\operatorname{seg}(w_0)\) onto \(\theta\), and \(u\prec v\) gives
\(F(u)\in F(v)\). It is injective, since distinct \(u,v\) are
\(\prec\)-comparable and then \(F(u)\in F(v)\) or \(F(v)\in F(u)\), each of
which excludes \(F(u)=F(v)\) by item~3 of
\cref{thm:p01-w02-p01-011-ordinals}. It reflects the order: if
\(F(u)\in F(v)\), then \(u\ne v\), and \(v\prec u\) would give
\(F(v)\in F(u)\), which item~3 excludes; so \(u\prec v\). Hence
\(\operatorname{seg}(w_0)\) is order isomorphic to \(\theta\), and
\(w_0\in D\), a contradiction. So \(D=W\), and \(F(w)\) is defined and
unique for every \(w\in W\) by item~1.

If \(v\prec w\), let \(h_w:\operatorname{seg}(w)\to F(w)\) be an order
isomorphism; the restriction step and uniqueness give
\(F(v)=h_w(v)\in F(w)\). Injectivity of \(F\) follows from comparability as
in the existence argument.

\emph{Item~3.} Let \(w\in W\) and let \(h:\operatorname{seg}(w)\to F(w)\) be
an order isomorphism. By item~3 of \cref{thm:p01-w02-p01-011-ordinals},
\(F(w)\in\kappa\), \(F(w)=\kappa\) or \(\kappa\in F(w)\). If
\(F(w)=\kappa\), then \(h\) is a bijection
\(\operatorname{seg}(w)\to\kappa\), against
\(\abs{\operatorname{seg}(w)}<\kappa\). If \(\kappa\in F(w)\), then
\(\kappa\subseteq F(w)\), so \(h^{-1}\) restricted to \(\kappa\) is an
injection \(\kappa\to\operatorname{seg}(w)\); with the injection
\(\operatorname{seg}(w)\to\kappa\) given by the hypothesis, CSB yields a
bijection \(\operatorname{seg}(w)\to\kappa\), again against the hypothesis.
So \(F(w)\in\kappa\). By item~2, \(F\) is an injection \(W\to\kappa\).
\end{proof}
\end{annexitem}

\begin{annexitem}{annex:p01-w02-p01-011-square}{thm:p01-w02-p01-011-square}{The square of an infinite cardinal}
\begin{proof}
Suppose some infinite cardinal \(\mu_0\) satisfies
\(\abs{\mu_0\times\mu_0}\ne\abs{\mu_0}\). By separation, the set \(S\) of
infinite cardinals \(\mu\in\mu_0+1\) with \(\abs{\mu\times\mu}\ne\abs\mu\)
exists, and \(\mu_0\in S\). Let \(\mu\) be its least member (item~4 of
\cref{thm:p01-w02-p01-011-ordinals}). If \(\lambda\) is an infinite cardinal
with \(\lambda<\mu\), then \(\lambda\in\mu\subseteq\mu_0+1\), because
\(\mu\in\mu_0+1\) and \(\mu_0+1\) is transitive; as \(\lambda\) lies below
the least member of \(S\), \(\lambda\notin S\), so
\(\abs{\lambda\times\lambda}=\abs\lambda\).

Order \(\mu\times\mu\) by \(\lhd\), a strict well-order by item~1 of
\cref{thm:p01-w02-p01-011-canonical-order}. We show
\(\abs{\operatorname{seg}(p)}<\mu\) for every \(p=(\alpha,\beta)\in\mu\times\mu\).
Put \(\gamma=\max(\alpha,\beta)\), a member of \(\mu\). By item~3 of
\cref{thm:p01-w02-p01-011-infinite-ordinals}, \(\gamma+1<\mu\), and by
item~2 of \cref{thm:p01-w02-p01-011-canonical-order},
\(\operatorname{seg}(p)\subseteq(\gamma+1)\times(\gamma+1)\).

\emph{Case 1: \(\gamma+1\) is finite.} Then \((\gamma+1)\times(\gamma+1)\) is
finite by item~5 of \cref{thm:p01-w02-p01-011-finite-sets}, its subset
\(\operatorname{seg}(p)\) is finite by item~3 there, and item~6 of
\cref{thm:p01-w02-p01-011-infinite-ordinals} gives
\(\abs{\operatorname{seg}(p)}<\mu\).

\emph{Case 2: \(\gamma+1\) is infinite.} Let \(\lambda=\operatorname{card}(\gamma+1)\) as in
item~4 of \cref{thm:p01-w02-p01-011-infinite-ordinals}: an infinite cardinal
with \(\lambda\leq\gamma+1\). Since \(\gamma+1\in\mu\) and \(\mu\) is
transitive, \(\lambda\in\mu\). So \(\abs{\lambda\times\lambda}=\abs\lambda\)
by the first paragraph. Let \(e:\gamma+1\to\lambda\) be a bijection; then
\((\xi,\eta)\mapsto(e(\xi),e(\eta))\) is a bijection
\((\gamma+1)\times(\gamma+1)\to\lambda\times\lambda\). Composing the
inclusion of \(\operatorname{seg}(p)\), this bijection and a bijection
\(\lambda\times\lambda\to\lambda\) gives
\(\abs{\operatorname{seg}(p)}\leq\abs\lambda\) with \(\lambda<\mu\), and
item~5 of \cref{thm:p01-w02-p01-011-infinite-ordinals} gives
\(\abs{\operatorname{seg}(p)}<\mu\).

Item~3 of \cref{thm:p01-w02-p01-011-representation} now gives an injection
\(\mu\times\mu\to\mu\). The map \(\xi\mapsto(\xi,0)\) is an injection
\(\mu\to\mu\times\mu\), since \(0\in\mu\) by item~1 of
\cref{thm:p01-w02-p01-011-infinite-ordinals}. CSB gives
\(\abs{\mu\times\mu}=\abs\mu\), against \(\mu\in S\). Hence no such
\(\mu_0\) exists.

No step uses a choice principle: the sets are formed by separation and
replacement, least members exist by item~4 of
\cref{thm:p01-w02-p01-011-ordinals}, the functions \(F\) of
\cref{thm:p01-w02-p01-011-representation} are uniquely determined, and the
proof of CSB in \cref{sec:set-prerequisites-05} defines its bijection
explicitly. The bijections \(e\) and \(\lambda\times\lambda\to\lambda\) are
single existential instantiations inside the proof of one inequality.
\end{proof}
\end{annexitem}

\begin{annexitem}{annex:p01-w02-p01-011-finite-coding}{thm:p01-w02-p01-011-finite-coding}{Finite coding over an infinite cardinal}
\begin{proof}
Fix a bijection \(\pi:\kappa\times\kappa\to\kappa\), which exists by
\cref{thm:p01-w02-p01-011-square}. By item~1 of
\cref{thm:p01-w02-p01-011-infinite-ordinals}, \(\N\subseteq\kappa\).

\emph{Clause~1.} By induction on \(n\geq1\) we show that a bijection
\(\kappa^n\to\kappa\) exists. For \(n=1\), \(s\mapsto s(0)\) is a bijection
\(\kappa^1\to\kappa\): a function on \(1=\{0\}\) is determined by its value
at \(0\), and every \(\xi<\kappa\) is the value of \(\{(0,\xi)\}\). Let
\(b_n:\kappa^n\to\kappa\) be a bijection. The map
\(s\mapsto(s\mathbin{\upharpoonright}n,s(n))\) is a bijection
\(\kappa^{n+1}\to\kappa^n\times\kappa\): it is injective because
\(n+1=n\cup\{n\}\), so \(s\) is determined by its restriction to \(n\) and
its value at \(n\); and it is surjective because for \(t\in\kappa^n\) and
\(\xi<\kappa\), \(t\cup\{(n,\xi)\}\) is a function on \(n+1\), as
\(n\notin n\). The map \((t,\xi)\mapsto(b_n(t),\xi)\) is a bijection
\(\kappa^n\times\kappa\to\kappa\times\kappa\), and \(\pi\) is a bijection
\(\kappa\times\kappa\to\kappa\). The composite is a bijection
\(\kappa^{n+1}\to\kappa\).

\emph{Clause~2.} Let \(K=\kappa^{<\omega}=\bigcup_{n\in\N}\kappa^n\). For
\(n\geq1\) let \(I_n\) be the set of injections \(\kappa^n\to\kappa\); it is
nonempty by clause~1. The family \(\{I_n:n\in\N,\ n\geq1\}\) is a set by
replacement, and its members are nonempty. The Axiom of Choice gives a
function \(c\) on this family with \(c(I_n)\in I_n\); put
\(g_n=c(I_n)\), an injection \(\kappa^n\to\kappa\). Every \(s\in K\) lies in
\(\kappa^n\) for exactly one \(n\), its domain. Define \(E:K\to\kappa\) by
\(E(\epsilon)=\pi(0,0)\) for the empty sequence \(\epsilon\), the only member
of \(\kappa^0\), and \(E(s)=\pi(n,g_n(s))\) for \(s\in\kappa^n\) with
\(n\geq1\); here \(n\in\N\subseteq\kappa\). If \(E(s)=E(t)\), injectivity of
\(\pi\) gives equal pairs. A first coordinate \(0\) occurs only for
\(\epsilon\). Otherwise \(s,t\in\kappa^n\) for the same \(n\geq1\) and
\(g_n(s)=g_n(t)\), so \(s=t\). So \(E\) is an injection \(K\to\kappa\). The
map \(\xi\mapsto\{(0,\xi)\}\) is an injection \(\kappa\to\kappa^1\subseteq K\).
CSB gives \(\abs K=\kappa\).

\emph{Clause~3.} Let \(\alpha<\kappa\), let \(A_\beta\) be finite for
\(\beta<\alpha\), and put \(U=\bigcup_{\beta<\alpha}A_\beta\).

If \(\alpha\) is finite, let \(e:n\to\alpha\) be a bijection with
\(n\in\N\). Then \(U=\bigcup_{i<n}A_{e(i)}\) is finite by item~4 of
\cref{thm:p01-w02-p01-011-finite-sets}, and item~6 of
\cref{thm:p01-w02-p01-011-infinite-ordinals} gives \(\abs U<\kappa\). No
choice is used in this case.

Let \(\alpha\) be infinite. For \(\beta<\alpha\) let \(E_\beta\) be the set
of bijections from some \(n\in\N\) onto \(A_\beta\); it is a set of subsets
of \(\N\times A_\beta\), and it is nonempty because \(A_\beta\) is finite.
The Axiom of Choice applied to \(\{E_\beta:\beta<\alpha\}\) gives, for every
\(\beta<\alpha\), a bijection \(b_\beta:n_\beta\to A_\beta\). For
\(x\in U\), let \(\beta(x)\) be the least \(\beta<\alpha\) with
\(x\in A_\beta\) (item~4 of \cref{thm:p01-w02-p01-011-ordinals}), and put
\(j(x)=(\beta(x),b_{\beta(x)}^{-1}(x))\in\alpha\times\N\). If
\(j(x)=j(y)\), then \(\beta(x)=\beta(y)=\beta\) and
\(b_\beta^{-1}(x)=b_\beta^{-1}(y)\), so \(x=y\). Since \(\alpha\) is
infinite, \(\N\subseteq\alpha\) by item~1 of
\cref{thm:p01-w02-p01-011-infinite-ordinals}, so \(j\) is an injection
\(U\to\alpha\times\alpha\). Let \(\lambda=\operatorname{card}(\alpha)\) (item~4 there), an
infinite cardinal with \(\lambda\leq\alpha<\kappa\), hence
\(\lambda<\kappa\). A bijection \(\alpha\to\lambda\) induces a bijection
\(\alpha\times\alpha\to\lambda\times\lambda\), and
\cref{thm:p01-w02-p01-011-square} gives a bijection
\(\lambda\times\lambda\to\lambda\). So \(\abs U\leq\abs\lambda\), and item~5
of \cref{thm:p01-w02-p01-011-infinite-ordinals} gives \(\abs U<\kappa\).

\emph{Clause~4.} For \(n\in\N\) let \(J_n\) be the set of injections
\(A_n\to\kappa\); by hypothesis each \(J_n\) is nonempty. The Axiom of
Choice applied to the set \(\{J_n:n\in\N\}\) gives a function \(c\) with
\(c(J_n)\in J_n\); put \(f_n=c(J_n)\). This selects the whole sequence
\((f_n)_{n\in\N}\) of injections \(f_n:A_n\to\kappa\) at once; the hypothesis
alone gives each \(f_n\) separately. Let \(U=\bigcup_{n\in\N}A_n\). For
\(x\in U\), let \(n(x)\) be the least \(n\) with \(x\in A_n\); it exists by
item~4 of \cref{thm:p01-w02-p01-011-ordinals}, as natural numbers are
ordinals. Put \(j(x)=(n(x),f_{n(x)}(x))\), a member of
\(\N\times\kappa\subseteq\kappa\times\kappa\). If \(j(x)=j(y)\), then
\(n(x)=n(y)=n\) and \(f_n(x)=f_n(y)\), so \(x=y\). Hence \(\pi\circ j\) is an
injection \(U\to\kappa\), and \(\abs U\leq\kappa\).
\end{proof}
\end{annexitem}

\begin{annexitem}{annex:p01-w02-p01-011-grammar-strings}{thm:p01-w02-p01-011-grammar-strings}{Generated expressions are finite strings}
\begin{proof}
\emph{Item~1.} By the definition of concatenation in
\cref{sec:strings-01}, \(uv\) is the function on \(m+n\) whose value at
\(i<m\) is \(u(i)\) and whose value at \(m\leq i<m+n\) is \(v(i-m)\). Both
values lie in \(\Gamma\), because \(u\) and \(v\) take values in
\(\Gamma\). So \(uv\in\Gamma^{m+n}\subseteq\Gamma^*\). For a template
\(u_0,\ldots,u_k\) and \(x_1,\ldots,x_k\in\Gamma^*\), the string
\(u_0x_1u_1\cdots x_ku_k\) is obtained from \(u_0\) by \(2k\) successive
concatenations with members of \(\Gamma^*\); by induction on the number of
concatenations performed, each intermediate string is in \(\Gamma^*\). So
the template rule is a function \((\Gamma^*)^k\to\Gamma^*\).

\emph{Item~2.} Let \((B,\mathcal F)\) be a finitary grammar over \(\Gamma\)
and \(G\) the set it generates. The set \(\Gamma^*\) contains \(B\), and it
is closed under every \(f\in\mathcal F\), because \(f\) maps
\((\Gamma^*)^{k_f}\) into \(\Gamma^*\). So \(\Gamma^*\) is one of the sets
\(Y\) whose intersection defines \(G\); in particular that family of sets is
nonempty, and \(G\subseteq\Gamma^*\). Equivalently, structural induction
(\cref{sec:strings-02}) for the property ``\(x\in\Gamma^*\)'' has its base
case from \(B\subseteq\Gamma^*\) and its constructor case from the fact that
every rule returns a member of \(\Gamma^*\). By
\cref{def:p01-w02-p01-011-readings}, \(\Gamma^*=\Gamma^{<\omega}\).

\emph{Item~3.} A function \(s:n\to\Gamma\) is a subset of
\(n\times\Gamma\subseteq n\times\Gamma'\) that relates each \(i<n\) to
exactly one value, so it is also a function \(n\to\Gamma'\).
\end{proof}
\end{annexitem}

\begin{annexitem}{annex:p01-w02-p01-011-syntax-counting}{thm:p01-w02-p01-011-syntax-counting}{Syntax counting}
\begin{proof}
\emph{Item~1.} Let \(f:A\to\kappa\) be an injection; such an injection is the
meaning of \(\abs A\leq\kappa\). For \(s\in A^n\), \(f\circ s\in\kappa^n\). The
map \(s\mapsto f\circ s\) from \(A^{<\omega}\) to \(\kappa^{<\omega}\)
preserves length and is injective: if \(f\circ s=f\circ t\), the two have
the same domain \(n\), so \(s,t\in A^n\), and \(f(s(i))=f(t(i))\) gives
\(s(i)=t(i)\) for every \(i<n\). Clause~2 of
\cref{thm:p01-w02-p01-011-finite-coding} gives a bijection
\(\kappa^{<\omega}\to\kappa\). The composite is an injection
\(A^{<\omega}\to\kappa\).

\emph{Item~2, strings.} The finite strings over \(\Sigma\) form
\(\Sigma^*=\Sigma^{<\omega}\). Item~1 with \(A=\Sigma\) gives
\(\abs{\Sigma^{<\omega}}\leq\kappa\).

\emph{Item~2, tuples.} A finite tuple of strings over \(\Sigma\) is a finite
sequence from \(\Sigma^{<\omega}\), so the tuples form
\((\Sigma^{<\omega})^{<\omega}\). By the previous step
\(\abs{\Sigma^{<\omega}}\leq\kappa\), and item~1 with
\(A=\Sigma^{<\omega}\) gives \(\abs{(\Sigma^{<\omega})^{<\omega}}\leq\kappa\).

\emph{Item~2, formulas.} Let the fixed finitary grammar have alphabet
\(\Gamma\supseteq\Sigma\), and let \(X=\Gamma\setminus\Sigma\) be the set of
adjoined symbols (logical symbols, punctuation and variables, for
instance). By hypothesis \(X\) is at most countable: there is an injection
\(X\to\N\), finite \(X\) included (\cref{sec:set-prerequisites-05}).
\begin{enumerate}
  \item \(\abs X\leq\kappa\). By item~1 of
        \cref{thm:p01-w02-p01-011-infinite-ordinals},
        \(\N\subseteq\kappa\), so the injection \(X\to\N\) followed by the
        inclusion is an injection \(X\to\kappa\). This is the inequality
        \(\abs X\leq\aleph_0\leq\kappa\).
  \item \(\abs\Gamma\leq\kappa\). Put \(A_0=\Sigma\), \(A_1=X\) and
        \(A_n=\varnothing\) for \(n\geq2\). Each \(A_n\) satisfies
        \(\abs{A_n}\leq\kappa\): \(A_0\) by hypothesis, \(A_1\) by step~1,
        and \(\varnothing\) through the empty function. Clause~4 of
        \cref{thm:p01-w02-p01-011-finite-coding} gives
        \(\abs{\bigcup_{n\in\N}A_n}\leq\kappa\), and
        \(\bigcup_{n\in\N}A_n=\Sigma\cup X=\Gamma\).
  \item The formula set is a set of finite strings over \(\Gamma\). The
        formulas are the expressions \(G\) generated by the grammar, and
        item~2 of \cref{thm:p01-w02-p01-011-grammar-strings} gives
        \(G\subseteq\Gamma^*=\Gamma^{<\omega}\).
  \item \(\abs G\leq\kappa\). Item~1 with \(A=\Gamma\), legitimate by
        step~2, gives an injection \(\Gamma^{<\omega}\to\kappa\); its
        restriction to \(G\) is an injection \(G\to\kappa\).
\end{enumerate}
The three sets of item~2 therefore each have cardinality at most
\(\kappa\).
\end{proof}
\end{annexitem}
